\chapter{Secure Web Application Design}
\label{ch:5.5}
\chaptertopic{Topic 5.5}

\chapterquote{There are two ways of constructing a software design: one way is to make it so simple that there are obviously no deficiencies, and the other way is to make it so complicated that there are no obvious deficiencies.}{C.~A.~R. Hoare, 1980 ACM Turing Award Lecture}

\section{Motivation}

Chapter~\ref{ch:5.1} showed what SQL injection, XSS, and BOLA look like from the attacker's side.
This chapter is the defender's answer, built as a discipline rather than a single fix. Three habits
make up that discipline: treat input as hostile by default, keep a codebase secure by construction
rather than by after-the-fact patching, and let automated testing find what those habits still
missed before an attacker does.

\begin{apalignment}{Unit 5: Securing Applications and Data}{Topic 5.5 -- Protecting Applications}
\begin{apobjectives}
\apobjective{5.5.A}{Mitigate Risk}{Explain the secure-coding habits (least privilege in code, parameterized queries, secrets management, failing securely, dependency management) that prevent vulnerabilities by construction, and how automated testing (SAST, DAST) plus human penetration testing verify those habits actually held.}
\apobjective{5.5.B}{Mitigate Risk}{Explain what checking untrusted input actually protects -- which interpreter is prevented from receiving syntax it would have obeyed -- and where the check has to run, why an allowlist beats a denylist, and which downstream defenses validation cannot replace.}
\end{apobjectives}
\end{apalignment}

\section{Input Validation and Sanitization}

Every application is fundamentally a machine that takes input and does something with it. Some
fraction of everyone typing into a form is not filling out their name -- they are trying to break the
application.

\begin{definition}{Input validation}
Checking that input conforms to what an application expects before it is used, as distinct from
\textbf{sanitization}, which transforms accepted input to remove or neutralize dangerous content
rather than rejecting it outright.
\end{definition}

Two philosophies govern validation. A \textbf{denylist} tries to enumerate the bad things --
reject any input containing \texttt{SELECT}, a \texttt{<script>} tag, a semicolon. It fails
constantly, because attackers are endlessly creative at encoding or rephrasing an attack past a
specific blocked pattern (\texttt{SELECT} becomes \texttt{SeLeCt}; a semicolon becomes its
URL-encoded \texttt{\%3B}). An \textbf{allowlist} flips the logic entirely: define exactly what
valid input looks like -- a ZIP code is five digits, a quantity is an integer between 1 and 99 --
and reject anything that does not match. Allowlisting is dramatically more robust because it never
has to anticipate every possible attack. It only has to know what correct looks like, and
correctness is almost always a far smaller, better-defined space than every way to attack a field.

\begin{lstlisting}[language=Python, caption={Allowlist validation for two field types}]
import re

FIELD_PATTERNS = {
    "zip_code": re.compile(r"^\d{5}$"),
    "email": re.compile(r"^[\w.+-]+@[\w.-]+\.[A-Za-z]{2,}$"),
}

def validate_field(field_type, value):
    """Returns True if value matches the allowlist pattern for field_type."""
    pattern = FIELD_PATTERNS.get(field_type)
    return pattern is not None and bool(pattern.match(value))
\end{lstlisting}

Validation performed only in the browser is worthless for security, because the browser is entirely
under the attacker's control. JavaScript can be disabled, the page can be edited with developer
tools, or a crafted HTTP request can bypass the browser entirely with a tool like curl. The rule to
memorize is unconditional -- never trust client-side validation for security; it must be re-enforced
on the server on every request, regardless of what the client claims it already checked.

Before validating at all, an application often must \textbf{canonicalize} input to a single,
predictable form. \texttt{../etc/passwd} and its URL-encoded form \texttt{..\%2Fetc\%2Fpasswd} are
the same logical string in two different representations. A validator that checks the raw,
un-canonicalized value can be tricked into approving something malicious, because it never recognized
what the string actually decodes to. The correct order is decode and normalize first, then validate
the canonical form. Validating before decoding is a classic bypass technique that smuggles
path-traversal or injection payloads directly past a naive filter.

\begin{misconception}
\textbf{``Sanitizing input once when it's received protects it everywhere it's later used.''} The
correct encoding depends entirely on the output context -- HTML, SQL, a shell command, JSON. A
single upstream sanitization pass is not a substitute for context-specific \textbf{output encoding}
applied at each point data is actually used.
\end{misconception}

SQL injection, XSS, and command injection all share the identical root failure: an application
trusted that input was ``just data.'' It actually contained characters a downstream interpreter --
a SQL engine, a browser, a shell -- treats as code. Rigorous input validation closes exactly this
gap at its source, before any of Chapter~\ref{ch:5.1}'s three attacks ever gets the chance to run.

\section{Secure Coding Practices}

Input validation stops one specific failure. \textbf{Secure coding} is the larger discipline: habits
and defaults built into every line of code so vulnerabilities are prevented by construction rather
than caught after the fact.

\textbf{Least privilege in code} extends Chapter~\ref{ch:4.2}'s principle beyond human accounts to
the code itself. A database connection should use an account that can only read and write the
specific tables it needs, never one with administrative rights over the entire database. A web
server process should never run as root, so that a single exploited bug becomes a contained failure
rather than a full system compromise. \textbf{Parameterized queries} are the single highest-leverage
habit against SQL injection specifically. Compare
\texttt{"SELECT * FROM users WHERE name = '" + name + "'"} -- dangerous, since the attacker's input
becomes query syntax -- to a query with \texttt{name} passed as a separate bound parameter, which is
safe because the driver treats it strictly as data no matter what characters it contains. One habit,
applied everywhere a query touches user input, eliminates the risk regardless of what validation
happened upstream. That is exactly why it counts as defense in depth alongside validation, rather
than a replacement for it.

\begin{lstlisting}[language=Python, caption={Auditing an application configuration against secure-coding practices}]
def audit_secure_coding(config):
    """config: dict like {"uses_parameterized_queries": bool,
       "hardcoded_secrets": bool, "runs_as_root": bool,
       "verbose_errors_in_production": bool}.
       Returns sorted violated-practice labels."""
    violations = []
    if not config.get("uses_parameterized_queries", False):
        violations.append("string-built-sql")
    if config.get("hardcoded_secrets", False):
        violations.append("hardcoded-secret")
    if config.get("runs_as_root", False):
        violations.append("excessive-process-privilege")
    if config.get("verbose_errors_in_production", False):
        violations.append("information-disclosure")
    return sorted(violations)
\end{lstlisting}

Secrets do not belong in source code at all: a hardcoded password or API key lives forever in
version-control history, visible to anyone with repository read access. It is trivially found by
both automated scanners and human review. Loading secrets at runtime from a secrets manager or
environment variable means rotating a compromised one requires no code change or redeployment.
\textbf{Failing securely} means an unexpected exception in an authentication check must deny access
by default, never accidentally grant it. What a failure \emph{reveals} matters too: a verbose error
page showing a stack trace or internal file path hands an attacker reconnaissance information, so
production code should log details internally while showing only a generic message externally.
Finally, \textbf{dependency management} recognizes that a vulnerability in any one of an
application's dozens or hundreds of third-party libraries becomes a vulnerability in the whole
application, exactly as Chapter~\ref{ch:4.3}'s patch-gap reasoning predicts. Excellent original code
sitting on a five-year-old unpatched library is not secure code.

\begin{misconception}
\textbf{``Secure coding is just input validation done well.''} Input validation is one practice
among several; parameterized queries, secrets management, least privilege, and dependency management
each close an entirely independent class of vulnerability that validation alone cannot touch.
\end{misconception}

\section{Application Security Testing}

Writing secure code is not the same as \emph{knowing} your code is secure. Two complementary
automated testing families answer that question from opposite directions.

\begin{definition}{SAST (static application security testing)}
Analyzing source code without executing it, searching for known-dangerous patterns -- a
string-concatenated SQL query, a call to \texttt{eval()} on user-controllable input, a hardcoded
secret. SAST catches these the moment code is written, often directly inside an editor or as a
CI/CD gate before a pull request merges.
\end{definition}

Because SAST never runs the program, it sees every code path, including rarely exercised error
handlers a live test might never stumble into. It catches vulnerabilities at the cheapest point
to fix, before deployment. Its cost is a high false-positive rate: a static tool cannot always tell
whether a flagged string concatenation actually reaches a database call with attacker-controlled
input, so it tends to over-flag. It also cannot see anything that only exists because of how the
deployed system is configured.

\begin{definition}{DAST (dynamic application security testing)}
Testing a running application from the outside as a black box -- sending real HTTP requests, fuzzing
form fields, manipulating cookies. It observes the responses for signs of a vulnerability, the
way an actual attacker would.
\end{definition}

DAST naturally catches runtime and configuration issues source code alone cannot reveal, and its
findings carry a lower false-positive rate. If the scanner injected a payload and got a genuinely
suspicious response back, something real is usually happening. Its cost is the mirror image of
SAST's -- it needs a running, reachable instance of the application, so it necessarily happens later
in the development cycle. It can only find what it actually tests: code paths its crawler never
reaches stay invisible.

\begin{lstlisting}[language=Python, caption={A miniature static-analysis scanner for dangerous code patterns}]
import re

DANGEROUS_PATTERNS = {
    "eval_call": re.compile(r"\beval\("),
    "string_built_sql": re.compile(r"(?i)select.*\+"),
    "hardcoded_secret": re.compile(r"(?i)(password|api_key)\s*=\s*[\"']"),
}

def scan_source_lines(lines):
    """Returns sorted (line_number, pattern_name) pairs for matches."""
    findings = []
    for i, line in enumerate(lines, start=1):
        for name, pattern in DANGEROUS_PATTERNS.items():
            if pattern.search(line):
                findings.append((i, name))
    return sorted(findings)
\end{lstlisting}

\begin{example}[Why the two approaches complement rather than compete]
SAST is cheap, early, and exhaustive over source code but blind to runtime behavior and prone to
false alarms. DAST is realistic and low-noise but late, incomplete in coverage, and blind to the
code that produced whatever it observes. A mature CI/CD pipeline runs SAST on every commit for fast
developer feedback and DAST against a staging deployment before release as a realistic pre-flight
check, then correlates the two. A finding confirmed by both tools is high-confidence. A SAST-only
finding needs a human to check reachability, and a DAST-only finding needs a human to trace back to
the responsible line of code.
\end{example}

\begin{figure}[h]
\centering
\resizebox{\textwidth}{!}{%
\begin{tikzpicture}[
  box/.style={draw=kickstartermaroon, thick, rounded corners, minimum width=2.6cm, minimum height=0.9cm, align=center, font=\scriptsize},
]
\node[box, fill=kickstartercream] (commit) at (0,0) {Every commit};
\node[box, fill=kickstartermaroon!15] (sast) at (3.2,0) {SAST\\(source, no run)};
\node[box, fill=kickstartercream] (staging) at (6.4,0) {Staging\\deployment};
\node[box, fill=kickstartermaroon!15] (dast) at (9.6,0) {DAST\\(live requests)};
\node[box, fill=kickstartermaroon!30] (release) at (12.8,0) {Release};
\draw[-{Latex[length=2mm]}] (commit) -- (sast);
\draw[-{Latex[length=2mm]}] (sast) -- (staging);
\draw[-{Latex[length=2mm]}] (staging) -- (dast);
\draw[-{Latex[length=2mm]}] (dast) -- (release);
\node[font=\scriptsize\sffamily, gray, align=center] at (6.4,-1.1) {a finding confirmed by both is high-confidence; a finding from only one still needs a human};
\end{tikzpicture}%
}
\caption{SAST runs early and cheaply on every commit, before the application even exists as a
running system; DAST runs later against a real staging deployment. Neither stage substitutes for
the other.}
\label{fig:5.5-sast-dast}
\end{figure}

Neither replaces a human \textbf{penetration test}: a skilled person who chains together multiple
weak findings and reasons about business-logic flaws no automated tool can recognize. ``Can a
regular user manipulate this parameter to view another user's invoice?'' is exactly the BOLA
question from Chapter~\ref{ch:5.1}, and it takes human judgment, not a scanner, to notice it applies
here. Mature programs use SAST and DAST for continuous, high-volume coverage and reserve human
penetration testing for periodic, deeper assurance ahead of major releases.

\begin{misconception}
\textbf{``A clean SAST/DAST scan means the application has no vulnerabilities.''} Both tools only
find what they are specifically built to look for. Business-logic flaws and novel attack chains
typically require human penetration testing to uncover at all.
\end{misconception}

\section*{Review Questions}
\begin{reviewquestions}
\item Explain why an allowlist validator for a ZIP-code field is more robust against novel attack
variants than a denylist that blocks specific known-bad substrings.
\item An application validates a filename only in its raw, un-decoded form before checking for
path-traversal sequences. Explain how this ordering mistake lets an attacker smuggle a malicious
path past the filter, and identify the correct order of operations.
\item Explain why parameterized queries are described as defense in depth alongside input
validation, rather than a replacement for it.
\item A SAST scanner flags a string-concatenated SQL query as a potential vulnerability, but a human
reviewer determines the concatenated value never actually contains user input. Explain why this
illustrates SAST's characteristic tradeoff.
\item Explain why neither SAST nor DAST would likely catch a BOLA vulnerability where an
authenticated user can view another user's invoice by changing a URL parameter, and explain what
kind of testing would.
\item Using the CI/CD pipeline diagram, explain why running DAST only, and skipping SAST entirely,
would still leave a codebase without the cheapest and earliest opportunity this chapter describes
to catch a vulnerability, even if DAST eventually caught the same issue later.
\item A finding is confirmed independently by both a SAST scan on a commit and a DAST scan against
staging. Explain why this combination is treated as higher-confidence than either finding alone,
referencing each tool's characteristic weakness from this chapter.
\end{reviewquestions}

\begin{answerkey}
\item A denylist must anticipate every possible way an attacker might encode or rephrase a bad
value, which is an effectively unbounded space attackers keep finding new members of. An allowlist
instead only has to define what a valid five-digit ZIP code actually looks like -- a small,
well-defined space -- and rejects everything else by default. A novel attack variant it has never
seen still fails to match the pattern and is rejected automatically.
\answerlabel{5.5.B}
\item Checking the raw, un-decoded filename lets an attacker submit an encoded path-traversal
sequence, such as \texttt{..\%2Fetc\%2Fpasswd}, that does not visibly match the filter's raw-text
check for \texttt{../} but decodes to exactly that traversal once the application later normalizes
it. The correct order is to decode and canonicalize the input to its single predictable form first,
then validate that canonical form. That way the filter inspects the value the application will
actually act on.
\answerlabel{5.5.B}
\item Parameterized queries eliminate the SQL-injection risk at the database layer regardless of
what validation happened upstream, but they do nothing for XSS, path traversal, or command
injection in the same request. Input validation catches malformed or suspicious input broadly across
all of these classes before it is used anywhere. Each closes a different gap, so running both
together is layered defense in depth, not two attempts at the same fix.
\answerlabel{5.5.A}
\item This illustrates SAST's high false-positive rate. Because SAST never executes the program, it
cannot always determine whether a flagged pattern's actual runtime inputs are attacker-controlled,
so it tends to over-flag constructs that merely resemble a vulnerability. Human review is needed to
confirm reachability before treating a finding as real.
\answerlabel{5.5.A}
\item SAST only recognizes dangerous code \emph{patterns} in source, and correctly-written
authorization code that simply forgets to check object ownership contains no dangerous pattern to
flag. DAST sends automated requests and looks for anomalous responses, but a BOLA request returns a
normal-looking, well-formed response. It is only wrong because it belongs to someone else, a
business-logic judgment no automated tool is positioned to make. A human penetration test, which
reasons about business logic such as ``can this parameter be manipulated to view another user's
data,'' is the kind of testing that catches it.
\answerlabel{5.5.A}
\item SAST runs on source code before the application is ever deployed, catching a vulnerability at
the cheapest possible point to fix -- before a pull request even merges. DAST requires a running,
reachable instance of the application, so it necessarily happens later in the development cycle,
after code has already progressed further through the pipeline. In a real organization, more work
has potentially been built on top of the flawed code by then. Skipping SAST means giving up that
early, low-cost catch even if DAST eventually finds the identical issue. The fix now
competes with everything built on top of it in the meantime, and the cost of fixing a defect
generally grows the later it is caught.
\answerlabel{5.5.A}
\item SAST's characteristic weakness is a high false-positive rate, since it cannot always confirm
a flagged pattern is actually reachable with attacker-controlled input. DAST's findings carry a
lower false-positive rate but a narrower scope, since it can only find what it actually tests.
When both tools independently flag the same issue through their very different methods -- one
reading source, the other sending live requests -- the chance that both are wrong about the exact
same finding for unrelated reasons is far lower than either tool being wrong alone. That is why
convergent findings from both are treated as high-confidence.
\answerlabel{5.5.A}
\end{answerkey}
